\section{Turning the entire enclosing prism}
The base is Keller's planar universal quadrilateral $C$ \cite[Theorem 4.1]{keller}. For $n(t)=(1-t^2,2t)/(1+t^2)$, its definition is
\[
 C=\{v\in\mathbb R^2:n(t_i)\cdot v\le d_i,\quad 0\le i\le3\},\qquad
 \begin{array}{c|rrrr}
 i&0&1&2&3\\ \hline
 t_i&-631/58&-1&7/16&37/56\\
 d_i&0&0&130783/250000&446199/10^6
 \end{array}
\]
 Index its four support lines cyclically from
$0$ through $3$, and write $n_i=(a_i,b_i)$ and $d_i$ for their unit
normals and offsets. The origin
is a vertex. Denote the intersection of lines $j$ and $j+1$ by $(x_j,y_j)$,
using indices modulo four. Write $x_* = x_3=\min_jx_j$ and
$Y=y_3=\max_j y_j$. The two bottom vertices have $y_0=y_1=0$.
Set
\[
 T=1/Y,\qquad H=(9+T^2)^{-1/2}=0.281666408810\ldots.
\]
Begin with the tapered body
\[
 K_0=\{(x,y,z):0\le z\le H,\quad(x,y)\in\sqrt{1-9z^2}\,C\}.
\]
The scale is nonnegative and concave on $[0,H]$. Since $0\in C$,
its scaled copies are nested, so $K_0$ is convex. Add one point $E=(0,H,H)$ and take $K=\operatorname{conv}(K_0\cup\{E\})$.

For a nonplanar curve choose the normal to the alternating support planes of Ghomi \cite[Corollary 3.2]{ghomi}, put $h$ equal to their separation, and write
$p=\sqrt{1-9h^2}$. The curve meets one plane, then the other, then the first and then the other again. Each of the three intervening portions contributes at least $h$ to the total variation of the vertical coordinate. The triangle inequality for the integrated horizontal and absolute
vertical speeds bounds the projected length by $p$. Indeed, in arclength coordinates,
\[
  L_{xy}^2+\operatorname{TV}(z)^2=
  \left\|\int_0^1 (\|\gamma'_{xy}\|,|z'|)\,ds\right\|^2\le1 .
\] A planar curve can
be treated with $h=0$. By Keller's covering theorem and scaling, the whole
spatial curve fits a right prism with base $p C$ and height $h$. If the
planar placement uses a reflection, also reverse the vertical coordinate
and translate its range to $[0,h]$. This gives a proper spatial motion.
For $h\le H$, this prism fits directly at the bottom of $K_0$.

For $H\le h\le1/3$, turn the prism through a right angle and translate it.
Its image can be written
\[
 Q_h=\{(p(x-x_*),s,py):(x,y)\in C,\quad 0\le s\le h\}.
\]
For example, the rotation $(x,y,z)\mapsto(x,-z,y)$ followed by translations
gives the displayed set. Only the eight vertices need to be tested because
$K$ is convex. Put $t=p/h$, so $0\le t\le T$. Since $TY=1$ and $p$
decreases with $h$, the largest height $pY$ occurs at $H$ and is at
most $HTY\le H$. Both horizontal coordinates of $Q_h$ are nonnegative.
Thus the first two inequalities of the horizontal sections hold.
The remaining normals have positive coordinates, so only $s=h$ can be
extremal. For $i=2,3$ and $X_j=x_j-x_*$ the required inequalities are
\[
 a_i p X_j+b_i h\le d_i\sqrt{1-9p^2 y_j^2}.
\]
Squaring, dividing by $h^2$, and using $h^{-2}=9+t^2$ gives the test
$A t^2+B t+D\le0$, where
\[
 A=(a_i X_j)^2+9d_i^2y_j^2-d_i^2,\quad B=2a_iX_j b_i,\quad
 D=b_i^2-9d_i^2.
\]
For seven entries rational upper bounds for the maximum over $[0,1/Y]$
are given below. One evaluates the polynomial at $0$, at $1/Y$, and at
$-B/(2A)$ when $A<0$ and this last point belongs to the interval.
\[
\begin{array}{c|rrrr}
 &j=0&j=1&j=2&j=3\\\hline
 i=2&-1.91&-0.014&-0.60&-0.44\\
 i=3&-0.93&-0.055&-0.18&\text{extra point}
\end{array}
\]
Only the vertex $(0,h,pY)$ can escape $K_0$. Let $m=d_3/b_3$.
Then $(0,m,0)\in K_0$. At every height $0\le z\le H$ the point $(0,0,z)$
lies in $K_0$, and the segment from $(0,m,0)$ to $E$ lies in $K$.
The resulting allowable interval of $y$ at $x=0$ reaches
$m-(m-H)z/H$. Along the exceptional vertex locus,
$y=g(z)=\sqrt{1-z^2/Y^2}/3$, for $0\le z\le H$.
Here $g(H)=H$ and $|g'(z)|\le |g'(H)|=1/(9Y^2)\le(m-H)/H$.
The available interval contains this vertex as well.

The addition of $E$ changes just one curved face of $K_0$.
Indeed, $E$ satisfies the defining inequalities of the bottom, top and
first three side surfaces viewed separately. Each such inequality defines
a convex set, so none of its supporting planes has $E$ strictly outside.
Let $e$ be the base edge between vertices 2 and 3, $L$ its length,
$s(z)=\sqrt{1-9z^2}$, $b=b_3$ and $d=d_3$.
The volume added by $E$ is the integral of its positive signed distance to
the supporting plane, times surface area divided by three, over the fourth
face. This cone formula follows by joining each visible surface element to
$E$. Equivalently, the parametrization $(r,u)\mapsto E+r(u-E)$ of the added
region, with $0\le r\le1$ and visible boundary point $u$, has volume
element $r^2(E-u)\cdot n(u)\,dr\,dS(u)$, where $n(u)$ is the outward
unit normal. Shared boundary pieces have volume zero.

At height $z$ the edge has length $L s(z)$ and the outward normal is
proportional to $(n_3,-d s'(z))$. Normalization cancels the area
factor $\sqrt{1+d^2s'(z)^2}$. For $u$ on this edge,
the resulting integrand before dividing by three is
$L s(z)[bH-ds(z)-ds'(z)(H-z)]_+$.
The identity $s(z)^2+s(z)s'(z)(H-z)=1-9Hz$ yields
\[
 |K|=|C| H(1-3H^2)
  +\frac L3\int_0^H
  \left[bH\sqrt{1-9z^2}-d(1-9Hz)\right]_+\,dz
  =0.05295603908\ldots .
\]
This integral also has an elementary closed form. With $c=bH/d$, let $z_0$ be the root in $(0,H)$ of
$(81H^2+9c^2)z_0^2-18Hz_0+1-c^2=0$. The expression inside the positive part is strictly concave, negative at zero and positive at $H$. Its positive interval is therefore $(z_0,H]$. An antiderivative there is
\[
 F(z)=bH\left(\tfrac z2\sqrt{1-9z^2}+\tfrac16\arcsin(3z)\right)-d\left(z-\tfrac92 Hz^2\right).
\]
The volume correction is exactly $L(F(H)-F(z_0))/3$.

For a rational evaluation put $z=Hv$. The expression inside the positive part becomes
$g(v)=bH\sqrt{1-9H^2v^2}-d(1-9H^2v)$.
It is concave, with $|g''|=9bH^3(1-9H^2v^2)^{-3/2}<3$ throughout $[0,1]$, using $1-9H^2>1/4$. Divide $[0,1]$ into $8192$ equal intervals. On each one $g$ exceeds its chord by at most $3/(8\cdot8192^2)$, by the linear interpolation error bound. Taking positive parts preserves that error. The integral of the positive part of a chord is the area of a trapezoid or a triangle. The verifier rounds each square root down to a multiple of $10^{-15}$ for these chords and adds the rounding and interpolation errors to the upper integral. It obtains the rigorous enclosure
\[
        0.0529560390 < \vol(K) < 0.052956040 .
\]
